%================================================================================
%       Safety Critical Systems Club - Data Safety Initiative Working Group
%================================================================================
%                       DDDD    SSSS  IIIII  W   W   GGGG
%                       D   D  S        I    W   W  G   
%                       D   D   SSS     I    W W W  G  GG
%                       D   D      S    I    WW WW  G   G
%                       DDDD   SSSS   IIIII  W   W   GGG
%================================================================================
%               Data Safety Guidance Document - LaTeX Source File
%================================================================================
%
% Description:
%   Objectives section.
%
%================================================================================
\chapter{Objectives (Informative)} \label{bkm:objectives}

\dsiwgSectionQuote
  {I'm a bit of a freak for evidence-based analysis. I strongly believe in data.}
  {Gus O'Donnell}


\section{General}

The data safety assurance principles (Chapter \ref{bkm:principles}) provide the underpinning philosophy for this guidance. An implementation of the guidance would be based around these objectives. This section identifies the objectives to satisfy the principles.


\section{Principle 1}
\dsiwgTextBF{\dsiwgTextIT{\index{Safety Requirement!Data}\Glspl{data safety requirement} shall be defined to address the data contribution to system hazards.}}

This guidance uses a metric called \glsfirst{dsal} for this purpose.\index{Assurance Level!Data}\index{DSAL|see{Assurance Level, Data}}, described in more detail in \dsiwgExtRef{Chapter}{2}{Part2-bkm:activities} and \dsiwgExtRef{Chapter}{2}{Part2-bkm:guidance}. The objectives and outputs supporting this principle are shown in Table \ref{tab:p1_principles}.
\begin{longtable}{|C{\dsiwgColumnWidth{0.06}}|L{\dsiwgColumnWidth{0.39}}|L{\dsiwgColumnWidth{0.39}}|L{\dsiwgColumnWidth{0.16}}|}
	\caption{P1 Objectives and Outputs}\label{tab:p1_principles}
	\\\hline\TableHeadColourC{Id} & \TableHeadColour{Objectives} & \TableHeadColour{Outputs} & \TableHeadColour{Reference} \\\hline
	\endfirsthead
	\caption{P1 Objectives and Outputs (continued)}
	\\\hline\TableHeadColourC{Id} & \TableHeadColour{Objectives} & \TableHeadColour{Outputs} & \TableHeadColour{Reference} 	\endhead
	\multicolumn{4}{r}{\sl Continued on next page}
	\endfoot\endlastfoot
	\hline
	1.a &%
	  System context and intended use SHALL be established.&%
	  A description of the system and its intended use e.g. in a \gls{conops}. This SHOULD include an estimate of the level of data risks.  %
	  & \dsiwgExtRef{Section}{2}{Part2-bkm:establishcontext}%
	  \\\hline
	1.b & %
	  Legal and regulatory obligations and requirements SHALL be incorporated. &%
	  A record of all applicable laws and regulations. & %
	  \dsiwgExtRef{Section}{2}{Part2-bkm:establishcontext}%
	  \\\hline
	1.c & %
	  Applicable standards and relevant guidance SHALL be incorporated & %
	  A record of all applicable standards and relevant guidance. & %
	  \hl{???Don't seem to address this in the guidance}%
	  \\\hline
	1.d & % 
	  Existing data quality and data assurance requirements and	constraints SHALL be identified & %
	  A record of applicable existing quality and data assurance requirements and constraints.& %
	  \hl{???Don't seem to address this in the guidance}%
	  \\\hline
	1.e & %
	  Data safety artefacts SHALL be identified &%
	  A collection of data artefacts, described at an appropriate level of detail. & % 
	  \dsiwgExtRef{Section}{2}{Part2-bkm:establishcontext} %
	  \\\hline
	\multirow{2}{*}{1.f} & %
	  \multirow{2}{6cm}{Risks SHALL be identified and linked to data safety artefacts} & %
	  A list of risks, linked to data artefacts and data properties. & %
	  \dsiwgExtRef{Section}{2}{Part2-bkm:identifyrisks} \\\cline{3-4} %
	 & & %
		A description of the process used for risk identification & %
		\dsiwgExtRef{Section}{2}{Part2-bkm:identifyrisks} %
		\\\hline
	1.g & %
	  Data safety activities SHALL be defined and documented &%
	  A plan for covering all data safety relevant aspects needs to be created. &% 
	  \dsiwgExtRef{Chapter}{2}{Part2-bkm:activities} %
	  \\\hline
	1.h & %
	  Key data stakeholders SHALL be identified &%
	  A list of key stakeholders for data safety activities, including those affected by the data. & %
	  \dsiwgExtRef{Section}{2}{Part2-bkm:establishcontext} %
	  \\\hline
	1.i & %
	  Missing, hidden, and obscured data SHALL be considered for safety impact &%
	  A plan to address potential relevant missing (Dark Data) and obscured data (Dazzle Data).&%
	  \hl{???Don't seem to have an activity for this, though could reference chapter on dark and dazzle data.} %
	  \\\hline
	1.j & %
	  External interfaces SHALL be defined. This MAY include a list of data owners, linked to data artefacts.&%
	  An interface control plan. &%
	  \hl{???Don't seem to have an activity for this.}
	  \\\hline
	1.k & %
	  A complete list of data assurance requirements SHALL be compiled, based on all the legal and regulatory needs, external inputs, known risks, etc. & %
	  A complete list of data assurance requirements. &
	  \dsiwgExtRef{Section}{2}{Part2-bkm:establishcontext}%
	  \\\hline
	1.l & %
	  Relevant historical data-related accidents and incidents SHOULD be reviewed for additional requirements &%
	  A documented review of relevant historical data-related accidents and incidents. &%
	  \dsiwgExtRef{Chapter}{2}{Part2-bkm:guidance} %
	  \\\hline
	1.m & %
	  \glspl{dsal} SHALL be established and justified.&%
	  A list of \glspl{dsal} and their justifications.&%
	  \dsiwgExtRef{Chapter}{2}{Part2-bkm:guidance} %
	  \\\hline
\end{longtable}
\cbend

\section{Principle 2}
\dsiwgTextBF{\dsiwgTextIT{The intent of \index{Safety Requirement!Data}\glspl{data safety requirement} shall be maintained throughout requirements decomposition.}}

This principle is based on standard systems engineering practices whereby a system is gradually developed at increasing levels of design detail. In order to cater for a wide range of systems, across a wide range of economic sectors, \gls{dsg} does not specifically require an explicit hierarchical decomposition of requirements. However, it does note that \index{Artefact, Data}\glspl{data artefact} MAY be defined at a number of levels of increasing detail. The objectives and outputs supporting this principle are shown in Table \ref{tab:p2_principles}.

\begin{longtable}{|C{\dsiwgColumnWidth{0.06}}|L{\dsiwgColumnWidth{0.39}}|L{\dsiwgColumnWidth{0.39}}|L{\dsiwgColumnWidth{0.16}}|}
	\caption{P2 Objectives and Outputs}\label{tab:p2_principles}
	\\\hline\TableHeadColourC{Id} & \TableHeadColour{Objectives} & \TableHeadColour{Brief Explanation and Deliverables} & \TableHeadColour{Reference in the Guidance} 	\endfirsthead
	\caption{P2 Objectives and Outputs (continued)}
	\\\hline\TableHeadColourC{Id} & \TableHeadColour{Objectives} & \TableHeadColour{Brief Explanation and Deliverables} & \TableHeadColour{Reference in the Guidance} 	\endhead
	\multicolumn{4}{r}{\sl Continued on next page}
	\endfoot\endlastfoot
	\hline
	2.a & %
	  Levels of abstraction of the data SHALL be defined. &%
	  Levels of Abstraction report. This MAY be in diagrammatic form. There MAY be several levels of data  abstraction, presenting different relevant perspectives. &%
      \hl{???Don't seem to address this in the guidance}%
      \\\hline
	2.b &%
	  \glspl{dsal} SHALL be flowed through levels of the design &%
	  Traceability records showing the flow of \glspl{dsal} through the design.&%
	  \hl{???Don't seem to address this in the guidance}%
	  \\\hline
	2.c &%
	  Data assurance requirements SHALL be established for each level of system design. &%
	  Requirements documents. &%
	  \hl{???Don't seem to address this in the guidance}%
	  \\\hline
	2.d &%
	  Data assurance requirements SHALL be traced between each level of system design &%
	  Traceability report &%
	  \hl{???Don't seem to address this in the guidance}%
	  \\\hline
	2.e &%
	  Data-related interactions between components SHALL be considered.&%
	  Interactions report. &%
	  \hl{???Don't seem to address this in the guidance}%
	  \\\hline
	2.f &%
	  Data-related interactions with other systems SHALL be considered.&%
	  Interactions report.&%
	  \hl{???Don't seem to address this in the guidance}%
	  \\\hline
	2.g &%
	  Techniques and approaches for mitigating data risks SHALL be identified&%
	  Mitigation report.&%
	  \hl{???Don't seem to address this in the guidance}%
	  \\\hline
	2.h &%
	  Unintended interactions affecting data SHALL be identified.&%
	  Interactions report.&%
	  \hl{???Don't seem to address this in the guidance}%
	  \\\hline
	2.i &%
	  Techniques and tools used for data analysis SHALL be justified.&%
	  Technique and tool justification.&%
	  \hl{???Don't seem to address this in the guidance}%
	  \\\hline
\end{longtable}
\cbend
\section{Principle 3}
\dsiwgTextBF{\dsiwgTextIT{\index{Safety Requirement!Data}\Glspl{data safety requirement} shall be satisfied.}}

The objectives and outputs supporting this principle are shown in Table \ref{tab:p3_principles}.

\begin{longtable}{|C{\dsiwgColumnWidth{0.06}}|L{\dsiwgColumnWidth{0.39}}|L{\dsiwgColumnWidth{0.39}}|L{\dsiwgColumnWidth{0.16}}|}
	\caption{P3 Objectives and Outputs}\label{tab:p3_principles}
	\\\hline\TableHeadColourC{Id} & \TableHeadColour{Objectives} & \TableHeadColour{Brief Explanation and Deliverables} & \TableHeadColour{Reference in the Guidance}
	\endfirsthead
	\caption{P3 Objectives and Outputs (continued)}
	\\\hline\TableHeadColourC{Id} & \TableHeadColour{Objectives} & \TableHeadColour{Brief Explanation and Deliverables} & \TableHeadColour{Reference in the Guidance}
	\endhead
	\multicolumn{4}{r}{\sl Continued on next page}
	\endfoot\endlastfoot
	\hline
	3.a &%
	  Methods and techniques used to verify data safety requirements SHALL be documented and justified &%
	  A record of the treatment adopted for each of the identified risks &%
	  \hl{???Don't seem to address this in the guidance}%
	  \\\hline
	\multirow{2}{*}{3.b} &%
	  \multirow{2}{5cm}{Data safety requirements SHALL be verified.}&%
	  Evidence that the treatment has been successfully implemented E.g. through test, analysis, review or other methods&%
	  \hl{???Don't seem to address this in the guidance}%
	  \\\cline{3-4}
	  & & Verification reports &
	  \hl{???Don't seem to address this in the guidance}%
	  \\\hline
	3.c &%
	   Methods and tools used to provide supplemental and supporting assurance SHALL be documented.&%
	   & %
	   \hl{???Don't seem to address this in the guidance}%
	   \\\hline
	3.d &%
	  \Glspl{dsal} SHALL be satisfied. &%
	  &%
	  \hl{???Don't seem to address this in the guidance}%
	  \\\hline
	3.e &%
	  Trusted tools SHALL be used to produce verification evidence &%
	  \dsiwgExtRef{Chapter}{2}{Part2-bkm:tools}&%
	  \\\hline
	3.f &%
	  Verification evidence SHOULD include evidence acquired from field or historical usage where available &%
	  &%
	  \hl{???Don't seem to address this in the guidance}%
	  \\\hline
	\multirow{2}{*}{3.g} &%
	  \multirow{2}{5cm}{Conformance with data safety requirements SHALL be documented.}&%
	  Conformance matrix & %
	  \hl{???Don't seem to address this in the guidance}%
	  \\\cline{3-4}
	  & & An assessment as to whether the risk has been suitably mitigated (and, if not, plans for further mitigation activities). &%
	  \hl{???Don't seem to address this in the guidance}%
	  \\\hline
	3.h &%
	  Data safety requirements which are partially or not met SHALL identified and justified &%
	  &%
	  \hl{???Don't seem to address this in the guidance}%
	  \\\hline
	3.i &%
	  Data which is missing or hidden SHALL be assessed for impact. &%
	  A missing and hidden data resilience report.&%
	  \hl{???Don't seem to have an activity for this, though could reference chapter on dark and dazzle data.} %
	  \\\hline
\end{longtable}
\cbend

\section{Principle 4}
\dsiwgTextBF{\dsiwgTextIT{Hazardous system behaviour arising from the system's use of data shall be identified and mitigated\index{Mitigation}.}}

The objectives and outputs supporting this principle are shown in Table \ref{tab:p4_principles}.\hl{???consider whether it's appropriate to refer to Black swan data, etc, which are not normatively defined.}
\begin{longtable}{|C{\dsiwgColumnWidth{0.06}}|L{\dsiwgColumnWidth{0.39}}|L{\dsiwgColumnWidth{0.39}}|L{\dsiwgColumnWidth{0.16}}|}
	\caption{P4 Objectives and Outputs}\label{tab:p4_principles}
	\\\hline\TableHeadColourC{Id} & \TableHeadColour{Objectives} & \TableHeadColour{Brief Explanation and Deliverables} & \TableHeadColour{Reference in the Guidance}
	\endfirsthead
	\caption{P4 Objectives and Outputs (continued)}
	\\\hline\TableHeadColourC{Id} & \TableHeadColour{Objectives} & \TableHeadColour{Brief Explanation and Deliverables} & \TableHeadColour{Reference in the Guidance}
	\endhead
	\multicolumn{4}{r}{\sl Continued on next page}
	\endfoot\endlastfoot
	\hline
	4.a &%
	  Key properties of the data SHALL be analysed for failures which lead to safety risks.&%
	  Data-focussed Functional Failure Analysis. &%
	  \hl{???Don't seem to address this in the guidance}%
	  \\\hline
	4.b &%
	  Unintended behaviour resulting from data SHALL be identified and analysed. The analysis SHOULD include \gls{black swan data}\index{Cygnology!Black Swan Data}, \gls{dragon king data}\index{Cygnology!Dragon King Data}, \gls{perfect storm data}\index{Cygnology!Perfect Storm Data}, and \gls{Pudding Lane data}\index{Cygnology!Pudding Lane Data}.&%
	  Unintended Behaviours Report. &%
	  \hl{???Don't seem to address this in the guidance}%
	  \\\hline
	4.c &% Unintended behaviour resulting from
      reduced or absent data SHALL be identified and analysed. &%
      Unintended Behaviours Report.&%
	  \hl{???Don't seem to address this in the guidance}%
	  \\\hline
	4.d &%
	  Emergent and unanticipated effects due to data SHALL be investigated.&%
	  &%
	  \hl{???Don't seem to address this in the guidance}%
	  \\\hline
	4.e &%
	  Resilience measures to deal with unknowns of the data SHALL be considered. &%
	  &%
	  \hl{???Don't seem to address this in the guidance}%
	  \\\hline
	4.f &%
	  Unintended behaviour resulting from malicious data SHALL be identified and analysed.&%
	  &%
	  \hl{???Don't seem to address this in the guidance}%
	  \\\hline
\end{longtable}
\cbend


\section{Principle 5}
\dsiwgTextBF{\dsiwgTextIT{The confidence established in addressing the data safety assurance principles shall be commensurate to the contribution of data to system risk.}}

The objectives and outputs supporting this principle are shown in Table \ref{tab:p5_principles}.

\begin{longtable}{|C{\dsiwgColumnWidth{0.06}}|L{\dsiwgColumnWidth{0.39}}|L{\dsiwgColumnWidth{0.39}}|L{\dsiwgColumnWidth{0.16}}|}
	\caption{P5 Objectives and Outputs}\label{tab:p5_principles}
	\\\hline\TableHeadColourC{Id} & \TableHeadColour{Objectives} & \TableHeadColour{Brief Explanation and Deliverables} & \TableHeadColour{Reference in the Guidance}
	\endfirsthead
	\caption{P5 Objectives and Outputs (continued)}
	\\\hline\TableHeadColourC{Id} & \TableHeadColour{Objectives} & \TableHeadColour{Brief Explanation and Deliverables} & \TableHeadColour{Reference in the Guidance}
	\endhead
	\multicolumn{4}{r}{\sl Continued on next page}
	\endfoot\endlastfoot
	\hline
	5.a & %
	  \Glspl{dsal} SHALL be demonstrably applied.&%
	  &%
	  \hl{???Don't seem to address this in the guidance}%
	  \\\hline
	5.b &%
	  Assurance artefacts SHALL be produced with a level of rigour commensurate with the \gls{dsal} and data safety requirements.&%
	  &%
	  \hl{???Don't seem to address this in the guidance}%
	  \\\hline
	5.c &%
	  Activities, methods, analyses, and tools used to provide assurance SHALL be appropriate for the \gls{dsal} and stakeholder group. &%
	  &%
	  \hl{???Don't seem to address this in the guidance}%
	  \\\hline
	5.d &%
	  Specific confidence demands from stakeholder groups SHALL be addressed.&%
	  &%
	  \hl{???Don't seem to address this in the guidance}%
	  \\\hline
	5.e &%
	  Residual risks SHALL be documented. &%
	  &%
	  \hl{???Don't seem to address this in the guidance}%
	  \\\hline
  \end{longtable}
\cbend

\section{Principle 6}
\hl{We don't have a Principal 6!}
The objectives and outputs supporting this principle are shown in Table \ref{tab:p6_principles}.
\begin{longtable}{|C{\dsiwgColumnWidth{0.06}}|L{\dsiwgColumnWidth{0.39}}|L{\dsiwgColumnWidth{0.39}}|L{\dsiwgColumnWidth{0.16}}|}
	\caption{P6 Objectives and Outputs}\label{tab:p6_principles}
	\\\hline\TableHeadColourC{Id} & \TableHeadColour{Objectives} & \TableHeadColour{Brief Explanation and Deliverables} & \TableHeadColour{Reference in the Guidance}
	\endfirsthead
	\caption{P6 Objectives and Outputs (continued)}
	\\\hline\TableHeadColourC{Id} & \TableHeadColour{Objectives} & \TableHeadColour{Brief Explanation and Deliverables} & \TableHeadColour{Reference in the Guidance}
	\endhead
	\multicolumn{4}{r}{\sl Continued on next page}
	\endfoot\endlastfoot
	\hline
	6.a &%
      All changes that impact these principles and objectives SHALL be assessed and managed.&%
	  &%
	  \\\hline
    6.b &%
      Assurance artefacts SHALL be maintained.&%
      &%
      \\\hline
    6.c &%
      Measures SHALL be in place to monitor the behaviour of the SCS at appropriate intervals throughout its life.&%
      &%
      \\\hline
    6.d &%
      Use SHALL be monitored for change and an assurance impact analysis undertaken  where necessary.&%
      &%
      \\\hline
    6.e &%
      Evolution and incremental change SHALL be considered.&%
      &%
      \\\hline
    6.f &%
      Longevity and obsolescence of data SCS SHALL be considered.&%
      &%
      \\\hline
    6.g &%
      Lessons learned from successful operation, failures and incidents SHOULD be reviewed&%
      &%
      \\\hline
\end{longtable}

